% !TEX root = ../lectures.tex

{\large\bfseries Calculus 2}

\section{Lecture 1}
\subsection*{Series (Prelim)}

\begin{definition}\label{def:1.1}
A function of the form
\[
 f:\N\to\R,\qquad n\longmapsto a_n,
\]
where $\N$ is the set of all positive integers and $\R$ is the set of real numbers, is called a \emph{sequence} (of real numbers).
\end{definition}

\noindent\textbf{Notations.}
\[
 (a_n)_{n\in\N}=(a_n)_{n=1}^{\infty}=(a_n)=(a_1,a_2,\ldots,a_n,\ldots).
\]
For example,
\[
 (n^2)_{n\in\N}=(n^2)_{n=1}^{\infty}=(n^2)
 =(1^2,2^2,3^2,\ldots,n^2,\ldots)
 =(1,4,9,\ldots,n^2,\ldots)
\]
is the sequence of squares.

\begin{definition}\label{def:1.2}
Let $(a_n)_{n=1}^{\infty}$ be a sequence. Then the $n$-th \emph{partial sum} (of $(a_n)_{n=1}^{\infty}$) is defined to be
\[
S_n=
\begin{cases}
 a_1, & n=1,\\[2mm]
 S_{n-1}+a_n, & n\ge 2.
\end{cases}
\]
\end{definition}

For example, if $\left(\frac1n\right)_{n=1}^{\infty}$ is a given sequence, then
\[
S_1=\frac11=1,
\]
\[
S_2=S_1+\frac12=1+\frac12=\frac32,
\]
\[
S_3=S_2+\frac13=\frac32+\frac13=\frac{11}{6},
\]
and, in general,
\[
S_n=1+\frac12+\cdots+\frac1n.
\]

\begin{definition}\label{def:1.3}
Let $(a_n)_{n=1}^{\infty}$ be a sequence. Then the \emph{series}
\[
\sum_{n=1}^{\infty}a_n
\]
is defined to be
\[
\sum_{n=1}^{\infty}a_n\overset{\mathrm{def}}{=}\lim_{n\to+\infty}S_n,
\]
where $S_n$ is the $n$-th partial sum of $(a_n)_{n=1}^{\infty}$.
\end{definition}

\noindent\textbf{Notation.} Instead of $\sum_{n=1}^{\infty}a_n$ one can use
\[
a_1+a_2+a_3+a_4+\cdots+a_n+\cdots,
\]
the expanded form.

\subsection*{Examples}
\begin{enumerate}[label=\arabic*)]
\item
\[
\sum_{n=1}^{\infty}\frac1n
=1+\frac12+\cdots+\frac1n+\cdots
\]
is the \emph{harmonic series}.

\item
\[
\sum_{n=1}^{\infty}\alpha\rho^{n-1}
=\alpha+\alpha\rho+\alpha\rho^2+\cdots+\alpha\rho^{n-1}+\cdots
\]
is the \emph{geometric series}.

\item Every real number $x\in[0,1)$ can be written as
\[
x=0.a_1a_2\ldots a_n\ldots,
\qquad a_n\in\{0,1,\ldots,9\}.
\]
This notation is a shorthand for the series
\[
x=\sum_{n=1}^{\infty}\frac{a_n}{10^n}.
\]
For instance,
\[
0.2282828\ldots
=\frac2{10}+\frac2{10^2}+\frac8{10^3}+\frac2{10^4}
 +\frac8{10^5}+\frac2{10^6}+\cdots.
\]

\item
\[
\sum_{n=1}^{\infty}(-1)^{n+1}
=1-1+1-1+\cdots+(-1)^{n+1}+\cdots.
\]
\end{enumerate}

\begin{definition}\label{def:1.4}
Let $(a_n)_{n=1}^{\infty}$ be a sequence. Then the series $\sum_{n=1}^{\infty}a_n$ is called
\begin{itemize}
\item \emph{convergent} (to $S$) if
\[
\lim_{n\to+\infty}S_n=S\in\R,
\qquad S_n=a_1+\cdots+a_n
\]
is the $n$-th partial sum;
\item \emph{divergent to $\pm\infty$} if
\[
\lim_{n\to+\infty}S_n=\pm\infty;
\]
\item \emph{divergent} if $\lim_{n\to+\infty}S_n$ does not exist.
\end{itemize}
\end{definition}

\subsection*{Examples of convergent series}
\begin{enumerate}[label=\arabic*)]
\item If $|\rho|<1$, then
\begin{align*}
\sum_{n=1}^{\infty}\alpha\rho^{n-1}
&\overset{\mathrm{def}}{=}
\lim_{n\to+\infty}\sum_{k=1}^{n}\alpha\rho^{k-1}\\
&=\lim_{n\to+\infty}\alpha\frac{1-\rho^n}{1-\rho}
=\alpha\frac{1-0}{1-\rho}
=\frac{\alpha}{1-\rho}\in\R.
\end{align*}

\item
\begin{align*}
\sum_{n=1}^{\infty}\frac1{n(n+1)}
&\overset{\mathrm{def}}{=}
\lim_{n\to+\infty}\sum_{k=1}^{n}\frac1{k(k+1)}\\
&=\lim_{n\to+\infty}\left(1-\frac1{n+1}\right)=1.
\end{align*}
This is a telescoping series.

\item
\[
\sum_{n=1}^{\infty}\frac1{n^2}
\overset{\mathrm{def}}{=}
\lim_{n\to+\infty}\sum_{k=1}^{n}\frac1{k^2}
=\frac{\pi^2}{6}.
\]
This is not trivial; we will prove it later in our course.

\item
\[
\sum_{n=1}^{\infty}\frac{(-1)^{n+1}}{n}
=\lim_{n\to+\infty}\sum_{k=1}^{n}\frac{(-1)^{k+1}}{k}
=\ln 2.
\]
\end{enumerate}

\subsection*{Examples of series divergent to $+\infty$}
\begin{enumerate}[label=\arabic*)]
\item
\[
\sum_{n=1}^{\infty}1
\overset{\mathrm{def}}{=}
\lim_{n\to+\infty}\sum_{k=1}^{n}1
=\lim_{n\to+\infty}n=+\infty.
\]

\item If $\rho>1$ and $\alpha>0$, then
\[
\sum_{n=1}^{\infty}\alpha\rho^{n-1}
\overset{\mathrm{def}}{=}
\lim_{n\to+\infty}\alpha\frac{1-\rho^n}{1-\rho}
=+\infty.
\]

\item
\[
\sum_{n=1}^{\infty}\frac1n
=\lim_{n\to+\infty}\sum_{k=1}^{n}\frac1k
=+\infty.
\]
We will prove this later.
\end{enumerate}

\subsection*{Examples of divergent series}
\begin{enumerate}[label=\arabic*)]
\item
\[
\sum_{n=1}^{\infty}(-1)^{n+1}
=\lim_{n\to+\infty}\sum_{k=1}^{n}(-1)^{k+1},
\]
where
\[
\sum_{k=1}^{n}(-1)^{k+1}
=
\begin{cases}
1, & n\text{ is odd},\\
0, & n\text{ is even}.
\end{cases}
\]
Hence the limit does not exist.

\item
\[
\sum_{n=1}^{\infty}\sin(n)
\]
is divergent: the limit $\lim_{n\to+\infty}\sum_{k=1}^{n}\sin(k)$ does not exist. We will prove this later.
\end{enumerate}

\begin{claim}\label{claim:1.1}
Changing a finite number of terms in a series does not affect its convergence (although it may change the sum of the series).
\end{claim}

\begin{proof}
Let $\sum_{n=1}^{\infty}a_n$ be the initial series and let $\sum_{n=1}^{\infty}b_n$ be obtained from it by changing a finite number of terms. Then it is clear that there exists $N\in\N$ such that
\[
\forall n\in\N,\quad n\ge N\implies a_n=b_n.
\]
Therefore, for every $n\ge N$,
\[
A_n-B_n=A_N-B_N,
\]
where
\[
A_k=\sum_{i=1}^{k}a_i,
\qquad
B_k=\sum_{i=1}^{k}b_i
\]
are the $k$-th partial sums. Thus
\[
\lim_{n\to+\infty}(A_n-B_n)=A_N-B_N\in\R,
\]
and the right-hand side does not depend on $n$. By the limit laws for sequences (see Calculus 1), it follows that either both limits $\lim A_n$ and $\lim B_n$ exist, or neither exists. Hence the convergence of the series is unchanged.
\end{proof}

\begin{claim}[Cauchy's criterion for series]\label{claim:1.2}
Let $\sum_{n=1}^{\infty}a_n$ be a series. Then
\[
\left(\sum_{n=1}^{\infty}a_n\text{ converges}\right)
\iff
\left(
\forall\varepsilon>0\ \exists N(\varepsilon)\in\N:\
\forall m,k\in\N,\ m\ge N(\varepsilon)
\implies
\left|\sum_{n=m+1}^{m+k}a_n\right|<\varepsilon
\right).
\]
\end{claim}

\begin{proof}
By \hyperref[def:1.4]{Definition~1.4},
\[
\sum_{n=1}^{\infty}a_n\text{ converges}
\iff
\exists\lim_{N\to+\infty}A_N,
\qquad
A_N=\sum_{n=1}^{N}a_n.
\]
By Cauchy's criterion for sequences (see Calculus 1), this is equivalent to
\[
\forall\varepsilon>0\ \exists N(\varepsilon)\in\N:\
\forall m,k\in\N,\ m\ge N(\varepsilon)
\implies |A_{m+k}-A_m|<\varepsilon.
\]
Since
\begin{align*}
A_{m+k}-A_m
&=(a_1+\cdots+a_m+a_{m+1}+\cdots+a_{m+k})-(a_1+\cdots+a_m)\\
&=a_{m+1}+\cdots+a_{m+k}
=\sum_{n=m+1}^{m+k}a_n,
\end{align*}
the claim is proved.
\end{proof}

\begin{claim}[The negation of the Cauchy criterion]\label{claim:1.3}
Let $\sum_{n=1}^{\infty}a_n$ be a series. Then
\[
\left(\sum_{n=1}^{\infty}a_n\text{ diverges}\right)
\iff
\left(
\exists\varepsilon>0:\ \forall N\in\N\ \exists m\ge N\ \exists k\in\N:\
\left|\sum_{n=m+1}^{m+k}a_n\right|\ge\varepsilon
\right).
\]
\end{claim}

\begin{proof}
Exercise (or one can see that it follows directly from \hyperref[claim:1.2]{Claim~1.2}).
\end{proof}

\begin{claim}\label{claim:1.4}
The following statements hold:
\begin{enumerate}[label=\arabic*)]
\item\label{item:claim:1.4:1} $\displaystyle \sum_{n=1}^{\infty}\frac1{n^2}$ converges;
\item\label{item:claim:1.4:2} $\displaystyle \sum_{n=1}^{\infty}\frac1n$ diverges.
\end{enumerate}
\end{claim}

\begin{proof}
\textbf{1)} To prove this fact we use \hyperref[claim:1.2]{Claim~1.2}. Let $\varepsilon>0$ be arbitrary. Set
\[
N(\varepsilon)=\left\lceil\frac1\varepsilon\right\rceil.
\]
For every $m\ge N(\varepsilon)$ and every $k\in\N$,
\begin{align*}
\left|\sum_{n=m+1}^{m+k}\frac1{n^2}\right|
&=\frac1{(m+1)^2}+\cdots+\frac1{(m+k)^2}\\
&\le \left(\frac1m-\frac1{m+1}\right)+\cdots+
\left(\frac1{m+k-1}-\frac1{m+k}\right)\\
&=\frac1m-\frac1{m+k}
<\frac1m
\le\frac1{N(\varepsilon)}
\le\varepsilon,
\end{align*}
where we used
\[
\frac1{n^2}\le \frac1{n(n-1)}=\frac1{n-1}-\frac1n.
\]
Therefore, by \hyperref[claim:1.2]{Claim~1.2}, $\sum_{n=1}^{\infty}\frac1{n^2}$ converges.

\medskip
\noindent\textbf{2)} We use \hyperref[claim:1.3]{Claim~1.3}. Let $\varepsilon=\frac12$. Let $N\in\N$ be arbitrary and set $m=N$, $k=N$. Then
\begin{align*}
\left|\sum_{n=m+1}^{m+k}\frac1n\right|
&=\sum_{n=N+1}^{2N}\frac1n\\
&=\frac1{N+1}+\frac1{N+2}+\cdots+\frac1{N+N}\\
&\ge \underbrace{\frac1{2N}+\frac1{2N}+\cdots+\frac1{2N}}_{N\text{ times}}
=N\cdot\frac1{2N}=\frac12=\varepsilon.
\end{align*}
Therefore, by \hyperref[claim:1.3]{Claim~1.3}, $\sum_{n=1}^{\infty}\frac1n$ diverges.
\end{proof}

\begin{claim}[Necessary condition for convergence]\label{claim:1.5}
Let $\sum_{n=1}^{\infty}a_n$ be a series. Then
\[
\left(\sum_{n=1}^{\infty}a_n\text{ converges}\right)
\implies
\left(\lim_{n\to+\infty}a_n=0\right).
\]
\end{claim}

\noindent\textbf{Note.} \hyperref[claim:1.4]{Claim~1.4}(2) shows that the converse does not hold. Indeed,
\[
\lim_{n\to+\infty}\frac1n=0,
\qquad
\sum_{n=1}^{\infty}\frac1n\text{ diverges}.
\]

\begin{proof}
This follows directly from \hyperref[claim:1.2]{Claim~1.2} if we set $k=1$.
\end{proof}

\begin{claim}\label{claim:1.6}
The series
\[
\sum_{n=1}^{\infty}\sin(nx)
\]
diverges for all $x\notin\pi\Z$.
\end{claim}

\begin{proof}
Suppose $\sum_{n=1}^{\infty}\sin(nx)$ converges. Then, due to \hyperref[claim:1.5]{Claim~1.5},
\[
\lim_{n\to+\infty}\sin(nx)
=\lim_{n\to+\infty}\sin((n+1)x)=0.
\]
Since
\[
\sin((n+1)x)=\sin(nx)\cos x+\cos(nx)\sin x
\]
and $\sin x\ne0$ (because $x\notin\pi\Z$), it follows that
\[
\lim_{n\to+\infty}\cos(nx)=0.
\]
But this contradicts the identity
\[
\sin^2(nx)+\cos^2(nx)=1.
\]
\end{proof}

\begin{definition}\label{def:1.5}
Let
\[
\sum_{n=1}^{\infty}a_n=S
\]
be a convergent series, and let
\[
S_n=\sum_{k=1}^{n}a_k
\]
be its $n$-th partial sum. Then, for every $n\in\N$,
\[
S=S_n+r_n,
\]
where the sequence $(r_n)_{n=1}^{\infty}$,
\[
r_n=\sum_{k=n+1}^{\infty}a_k,
\]
is called the \emph{tail} (or the \emph{remainder}) of the series $\sum_{n=1}^{\infty}a_n$.
\end{definition}

\begin{claim}\label{claim:1.7}
Let $\sum_{n=1}^{\infty}a_n$ be a series. Then
\[
\left(\sum_{n=1}^{\infty}a_n\text{ converges}\right)
\implies
\left(\lim_{n\to+\infty}r_n=0\right).
\]
\end{claim}

\begin{proof}
Let $\sum_{n=1}^{\infty}a_n=S\in\R$. Then
\[
r_n=S-S_n\qquad\forall n\in\N.
\]
Therefore,
\[
\lim_{n\to+\infty}r_n
=\lim_{n\to+\infty}(S-S_n)
=S-\lim_{n\to+\infty}S_n
=S-S=0.
\]
\end{proof}

\begin{claim}[Linear combination of series]\label{claim:1.8}
Let $\sum_{n=1}^{\infty}a_n$ and $\sum_{n=1}^{\infty}b_n$ be two convergent series. Then, for every $\alpha,\beta\in\R$, the series
\[
\sum_{n=1}^{\infty}(\alpha a_n+\beta b_n)
\]
is also convergent and
\[
\sum_{n=1}^{\infty}(\alpha a_n+\beta b_n)
=\alpha\sum_{n=1}^{\infty}a_n
+\beta\sum_{n=1}^{\infty}b_n.
\]
This series is the linear combination of the series $\sum a_n$ and $\sum b_n$ with coefficients $\alpha$ and $\beta$.
\end{claim}

\begin{proof}
Since $\sum_{n=1}^{\infty}a_n$ and $\sum_{n=1}^{\infty}b_n$ are convergent, for some $a,b\in\R$ we have
\[
\sum_{n=1}^{\infty}a_n=\lim_{n\to+\infty}A_n=a,
\qquad
A_n=\sum_{k=1}^{n}a_k,
\]
and
\[
\sum_{n=1}^{\infty}b_n=\lim_{n\to+\infty}B_n=b,
\qquad
B_n=\sum_{k=1}^{n}b_k.
\]
Therefore, by the limit laws for sequences (see Calculus 1),
\begin{align*}
\alpha a+\beta b
&=\alpha\lim_{n\to+\infty}A_n+\beta\lim_{n\to+\infty}B_n\\
&=\lim_{n\to+\infty}(\alpha A_n+\beta B_n)\\
&=\lim_{n\to+\infty}\left(\sum_{k=1}^{n}(\alpha a_k+\beta b_k)\right).
\end{align*}
By \hyperref[def:1.3]{Definition~1.3}, this is precisely
\[
\sum_{n=1}^{\infty}(\alpha a_n+\beta b_n)=\alpha a+\beta b.
\]
\end{proof}

\begin{claim}[Grouping theorem]\label{claim:1.9}
Let
\begin{enumerate}[label=\alph*)]
\item\label{item:claim:1.9:a} $\displaystyle \sum_{n=1}^{\infty}a_n$ be a series;
\item\label{item:claim:1.9:b} $(K_n)_{n=1}^{\infty}$ be a sequence of positive integers such that
\[
K_1=1,
\qquad
K_i<K_{i+1}\quad\text{for every }i\in\{1,2,\ldots\};
\]
\item\label{item:claim:1.9:c} $\displaystyle \sum_{n=1}^{\infty}b_n$, where
\[
b_n=a_{K_n}+\cdots+a_{K_{n+1}-1},
\qquad n\in\N,
\]
be a \emph{grouped series}.
\end{enumerate}
Thus $b_n$ is the $n$-th block and $K_{n+1}-K_n$ is the size of the $n$-th block.
Then:
\begin{enumerate}[label=\arabic*)]
\item\label{item:claim:1.9:1}
\[
\left(\sum_{n=1}^{\infty}a_n=\alpha\in\R\right)
\implies
\left(\sum_{n=1}^{\infty}b_n=\alpha\right).
\]
\item\label{item:claim:1.9:2} If
\[
\lim_{n\to+\infty}a_n=0,
\qquad
\exists m\in\N\ \forall n\in\N:\ K_{n+1}-K_n<m,
\qquad
\sum_{n=1}^{\infty}b_n=b\in\R,
\]
then
\[
\sum_{n=1}^{\infty}a_n=b.
\]
\end{enumerate}
\end{claim}
